% ============================================================================
% DÉFINITION DU COURS 03 : LOIS ENTRÉE-SORTIE
% Le contenu pédagogique reste dans la première passe contenu.tex.
% ============================================================================

\titlehead{Thomas Lusseau - SII en ATS}
\title[Lois entrée-sortie]{Détermination des lois de mouvement -- Lois entrée-sortie}
\author[TL]{Thomas Lusseau}
\date{\today}

\TLCourseCoverSetup
  {\repStyle/FondPoly.png}
  {images/image3.png}
  {Cycle 2 : Modéliser et déterminer les lois de\\ mouvement d'un mécanisme}
  {Détermination des lois de commande\\ Loi entrée-sortie}
  {Lycée Robert Doisneau - ATS}
  {Thomas Lusseau}

\TLCourseCoverContents{%
  \TLCourseContentsLine{3.1}{Mise en situation}{mise-en-situation}
  \TLCourseContentsLine{3.2}{Modèle géométrique des structures fermées}{modele-geometrique-de-structures-fermees}
  \TLCourseContentsLine{3.3}{Détermination d'une loi entrée-sortie}{determination-dune-loi-entree-sortie}
  \TLCourseContentsLine{3.4}{Lois entrée-sortie en vitesse}{lois-entree-sortie-en-vitesse}
  \TLCourseContentsLine{3.5}{Exercices du chapitre}{exercices-du-chapitre}
}

\TLCourseFooter
\TLCourseCover
\markboth{Lois entrée-sortie}{Lois entrée-sortie}

\AtBeginEnvironment{corrige}{\small}
\livrettrue
\collefalse

\setcounter{chapter}{2}
\chapter*{Détermination des lois de mouvement}
\addcontentsline{toc}{chapter}{Détermination des lois de mouvement}

% Partie cours de la première passe « Lois entrée-sortie ».
% Le fichier contenu.tex reste inchangé ; on n'exécute que ce qui précède
% la sous-section « Exercices du chapitre ».

\begingroup
% Couche locale d'intégration de la conversion Word du cours de lois E/S.
% Ce fichier doit être chargé dans un groupe TeX.

\makeatletter
\let\TLThreeOriginalIncludeGraphics\includegraphics

\RenewDocumentCommand{\includegraphics}{s O{} m}{%
  \begingroup
  \filename@parse{#3}%
  \edef\TLThreeGraphicBase{\filename@base}%

  % Pictogrammes injectés dans les titres par la conversion Word.
  \ifboolexpr{
    test {\ifdefstring{\TLThreeGraphicBase}{image8}}
    or test {\ifdefstring{\TLThreeGraphicBase}{image26}}
    or test {\ifdefstring{\TLThreeGraphicBase}{image30}}
    or test {\ifdefstring{\TLThreeGraphicBase}{image32}}
    or test {\ifdefstring{\TLThreeGraphicBase}{image33}}
    or test {\ifdefstring{\TLThreeGraphicBase}{image34}}
  }{%
    % rien
  }{%
    \ifdefstring{\TLThreeGraphicBase}{image22}{%
      \TLThreeOriginalIncludeGraphics[width=8mm,height=8mm,keepaspectratio]{#3}%
    }{%
      \ifdefstring{\TLThreeGraphicBase}{image11}{%
        \ifinner
          \TLThreeOriginalIncludeGraphics[width=.72\linewidth,height=.22\textheight,keepaspectratio]{#3}%
        \else
          \par\noindent
          \makebox[\linewidth][c]{%
            \TLThreeOriginalIncludeGraphics[width=.58\linewidth,height=.30\textheight,keepaspectratio]{#3}%
          }%
          \par\noindent
        \fi
      }{%
        \ifdefstring{\TLThreeGraphicBase}{image31}{%
          \ifinner
            \TLThreeOriginalIncludeGraphics[width=.78\linewidth,height=.30\textheight,keepaspectratio]{#3}%
          \else
            \par\noindent
            \makebox[\linewidth][c]{%
              \TLThreeOriginalIncludeGraphics[width=.72\linewidth,height=.34\textheight,keepaspectratio]{#3}%
            }%
            \par\noindent
          \fi
        }{%
          \ifinner
            \adjustbox{max width=.92\linewidth,max totalheight=.70\textheight}{%
              \IfBooleanTF{#1}
                {\TLThreeOriginalIncludeGraphics*[#2]{#3}}
                {\TLThreeOriginalIncludeGraphics[#2]{#3}}%
            }%
          \else
            \par\noindent
            \makebox[\linewidth][c]{%
              \adjustbox{max width=.94\linewidth,max totalheight=.76\textheight}{%
                \IfBooleanTF{#1}
                  {\TLThreeOriginalIncludeGraphics*[#2]{#3}}
                  {\TLThreeOriginalIncludeGraphics[#2]{#3}}%
              }%
            }%
            \par\noindent
          \fi
        }%
      }%
    }%
  }%
  \endgroup
}
\makeatother


\CatchFileDef{\TLThreeSource}{contenu.tex}{}

% Parcours récursif des sous-sections jusqu'au séparateur des exercices.
\long\def\TLThreeExtractCourse#1\subsection#2#3\TLThreeStop{%
  \ifstrequal{#2}{Exercices du chapitre}{%
    #1%
  }{%
    #1\subsection{#2}%
    \TLThreeExtractCourse#3\TLThreeStop
  }%
}

\expandafter\TLThreeExtractCourse
  \TLThreeSource
  \TLThreeStop
\endgroup

